% !TEX root = ../main.tex
% =========================================================
% bab5.tex - PENUTUP
% =========================================================

\chapter[PENUTUP]{\\ PENUTUP}

\section{Kesimpulan}

Berdasarkan analisis dan pembahasan data hasil pengujian kurva I--V modul PV
pasca-terendam pada PLTS Terapung Cirata, diperoleh beberapa kesimpulan, yaitu:

\begin{enumerate}
    \item Analisis dilakukan terhadap 59 data pengujian kurva I--V modul PV
          pasca-terendam tipe JKM560N-72HL4-BDV yang diuji secara
          \textit{onshore} dengan koreksi parameter ke STC. Karena terdapat
          satu nomor uji yang muncul dua kali, hasil dinyatakan berdasarkan
          data pengujian dan bukan sebagai klaim jumlah modul unik.

    \item Hasil \textit{screening} performa daya oleh perangkat lunak pengolah
          data menunjukkan 6 data (10,17\%) berstatus Ok dan 53 data (89,83\%)
          berstatus Not Ok. Status ini merupakan hasil perbandingan $P_{max}$
          STC terukur terhadap nilai \textit{nameplate} modul dan belum
          merepresentasikan keputusan akhir kelayakan operasional.

    \item Kelompok Not Ok mengalami penurunan rata-rata $P_{max}$ STC menjadi
          sekitar 501,6~W atau sekitar 10,4\% di bawah nilai \textit{nameplate}
          560~W. Penurunan ini menunjukkan bahwa modul pada kelompok Not Ok
          mengalami penurunan performa daya yang substansial terhadap
          spesifikasi pabrikan.

    \item Analisis deviasi indikator $R_{VM}$ dan $R_{IM}$ terhadap nilai
          referensi \textit{datasheet} menghasilkan $\Delta R_{VM} = -1{,}57\%$
          dan $\Delta R_{IM} = +1{,}71\%$ untuk kelompok Not Ok, sehingga
          diperoleh $\delta_V = 1{,}57\%$, $\delta_I = 0{,}00\%$, dan
          $D_V = 1{,}00$. Indikator $R_{VM}$ berbeda signifikan antara kedua
          kelompok berdasarkan Welch \textit{t-test} ($p < 0{,}001$), sedangkan
          $R_{IM}$ tidak berbeda signifikan ($p = 0{,}150$), sehingga pemisahan
          kelompok lebih dominan terjadi pada sisi tegangan.

    \item Pola penurunan sisi tegangan pada kelompok Not Ok konsisten dengan
          kemungkinan peningkatan resistansi seri atau mekanisme degradasi lain
          yang menurunkan tegangan operasi modul. Interpretasi ini dibatasi
          sebagai indikasi awal berbasis I--V dan tidak digunakan untuk
          menetapkan jenis kerusakan fisik modul secara final.

    \item Hasil pengujian I--V tidak cukup menjadi dasar tunggal keputusan
          pemasangan kembali karena aspek keselamatan listrik tidak dapat
          disimpulkan hanya dari parameter kurva I--V, sehingga diperlukan
          pengujian tahanan isolasi dan pemeriksaan keselamatan lanjutan.
\end{enumerate}

\section{Saran}

Berdasarkan hasil analisis yang telah dilakukan, berikut ini adalah saran-saran
yang dapat digunakan sebagai panduan untuk kegiatan dan pengembangan
selanjutnya:

\begin{enumerate}
    \item Melakukan pengujian tahanan isolasi (\textit{insulation resistance
          test}) pada seluruh modul pasca-terendam sebelum keputusan pemasangan
          kembali, karena risiko utama bukan hanya penurunan daya tetapi juga
          keselamatan listrik.

    \item Melengkapi evaluasi dengan inspeksi visual, termografi, dan/atau
          elektroluminesensi jika tersedia untuk mengonfirmasi mekanisme
          degradasi yang teridentifikasi dari parameter elektrik.

    \item Melakukan pengujian ulang pada kondisi irradiansi yang lebih
          representatif untuk data yang diperoleh pada irradiansi rendah,
          khususnya kelompok Ok, sebelum data tersebut dijadikan dasar
          keputusan.

    \item Mendokumentasikan kondisi modul dan riwayat lokasi modul secara
          sistematis agar interpretasi pola degradasi dapat dikaitkan dengan
          kondisi operasi aktual.

    \item Mengintegrasikan hasil pengujian I--V dengan prosedur
          \textit{operation and maintenance} PLTS Terapung Cirata, termasuk
          pemantauan parameter tegangan ($V_{mp}$ dan tegangan
          \textit{string}) sebagai indikator awal adanya modul terdegradasi.
\end{enumerate}